\documentclass[10pt,a4paper]{amsart}
\usepackage[margin=19mm]{geometry}
\usepackage{amsmath,amssymb,mathtools}
\usepackage[scr=rsfs,cal=euler]{mathalfa}
\usepackage{xfrac}
\usepackage[unicode,hidelinks,bookmarks=false]{hyperref}
\newcommand{\ie}{i.e.\ }
\newcommand{\cf}{cf.\ }
\newcommand{\resp}{respectively\ }
\newcommand{\Subsection}{\subsection}
\providecommand{\nobreakdash}{\nobreak\hbox{-}}
\setlength{\emergencystretch}{2em}
\multlinegap0pt
\usepackage{float}
\usepackage{amssymb}
\newtheorem{lem}{}
\title{On the Theory of Bias Tuning in Event Cameras}
\author{David El-Chai Ben-Ezra, Daniel Brisk, Adar Tal}
\date{Source: arXiv:2501.18788v3 (CC BY 4.0)}
\begin{document}
\maketitle

\noindent\emph{Source and licence.} On the Theory of Bias Tuning in Event Cameras. Version arXiv:2501.18788v3 (CC BY 4.0), distributed by arXiv under the Creative Commons Attribution 4.0 International licence, http://creativecommons.org/licenses/by/4.0/, as recorded in the arXiv licence metadata for this exact version and shown on the abstract page https://arxiv.org/abs/2501.18788v3. Full licence notice: \texttt{licenses/arxiv-2501.18788-CC-BY-4.0.txt}.\par\smallskip

\noindent\textit{Editorial scope.} Counted unit: the article's lem lem:2 (source0.tex lines 443--451). The article's complete original proof of it is delivered with it (source0.tex lines 453--497, one \texttt{proof} environment of 45 lines). No mathematics is written, repaired or re-derived: the statement and the proof are the article's own lines, in the article's own notation, and no retained line is substituted or re-flowed.  Retained uncounted as the source-local context the proof needs: lines 384--402.  Nothing was omitted from any retained span, so the package carries no omission record.  No retained line contains a citation, so the wrapper carries no bibliography and no bibliography entry.  No retained line contains a figure environment, an image-inclusion command or drawing code, so no image is delivered and the package declares no asset dependency.  Editorial formatting only, in addition: an AMS \texttt{amsart} preamble that restates the article's own declarations and macros used by the retained lines, namely the environments \texttt{lem} on separate counters, together with the standard packages those lines need.  The retained environments print in this document as Lemma 1 in order of appearance, whereas the article prints the same items with the numbering of its own full document.  The complete native source of the article as distributed in the arXiv e-print for this exact version is bundled unchanged as \texttt{sources/172486/source0.tex}.
\begin{lem}
\label{lem:main}Let $f(x,y)$ and $g(x,y)$ be real-valued differentiable
functions, defined over the rectangles $[a,b)\times[c,d]$ and $[a,b]\times[c,d)$,
correspondingly. Assume also that for every $(x_{0},y_{0})$ in the
regions we have:

\begin{align*}
1. & \underset{x\to b}{\lim}f(x,y_{0})=\underset{y\to d}{\lim}g(x_{0},y)=\infty\\
2. & \frac{\partial(\alpha\circ f)}{\partial x}(x_{0},y_{0})>\frac{\partial(\alpha\circ g)}{\partial x}(x_{0},y_{0})\geq0\\
3. & \frac{\partial(\alpha\circ g)}{\partial y}(x_{0},y_{0})>\frac{\partial(\alpha\circ f)}{\partial y}(x_{0},y_{0})\geq0
\end{align*}
for some monotonic $\alpha:\mathbb{R}^{+}\to\mathbb{R}$. Then, there
exist two constants $0\leq K_{1},K_{2}$ such that for every $K_{1}\leq k_{1}$
and $K_{2}\leq k_{2}$ there exists a unique solution to the equation
system
\[
f(x,y)=k_{1},\qquad g(x,y)=k_{2}.
\]
\end{lem}

\begin{lem}
\label{lem:2}Let $f$ and $g$ be as in Lemma \ref{lem:main}. Then,
both coordinates of the unique solution of the system
\[
f(x,y)=g(x,y)=k
\]
grow as functions of $k$. In other words, $x(k)$ and $y(k)$ are
monotonic.
\end{lem}

\begin{proof}
Let us denote $f_{\alpha}=\alpha\circ f$ and $g_{\alpha}=\alpha\circ g$,
and write the equations
\[
f_{\alpha}(x(k),y(k))=g_{\alpha}(x(k),y(k))=\alpha(k)
\]
where $x(k)$ and $y(k)$ are the unique values given by Lemma \ref{lem:main}
for each $k$. Then, by the chain rule
\[
\begin{array}{c}
\frac{\partial\alpha}{\partial k}=\frac{\partial f_{\alpha}}{\partial k}=\frac{\partial f_{\alpha}}{\partial x}\cdot\frac{\partial x}{\partial k}+\frac{\partial f_{\alpha}}{\partial y}\cdot\frac{\partial y}{\partial k}\\
\frac{\partial\alpha}{\partial k}=\frac{\partial g_{\alpha}}{\partial k}=\frac{\partial g_{\alpha}}{\partial x}\cdot\frac{\partial x}{\partial k}+\frac{\partial g_{\alpha}}{\partial y}\cdot\frac{\partial y}{\partial k}
\end{array}\Rightarrow\left(\begin{array}{cc}
\frac{\partial f_{\alpha}}{\partial x} & \frac{\partial f_{\alpha}}{\partial y}\\
\frac{\partial g_{\alpha}}{\partial x} & \frac{\partial g_{\alpha}}{\partial y}
\end{array}\right)\cdot\left(\begin{array}{c}
\frac{\partial x}{\partial k}\\
\frac{\partial y}{\partial k}
\end{array}\right)=\left(\begin{array}{c}
\frac{\partial\alpha}{\partial k}\\
\frac{\partial\alpha}{\partial k}
\end{array}\right)
\]
yielding that
\[
\left(\begin{array}{c}
\frac{\partial x}{\partial k}\\
\frac{\partial y}{\partial k}
\end{array}\right)=\frac{1}{\frac{\partial f_{\alpha}}{\partial x}\cdot\frac{\partial g_{\alpha}}{\partial y}-\frac{\partial f_{\alpha}}{\partial y}\cdot\frac{\partial g_{\alpha}}{\partial x}}\cdot\left(\begin{array}{cc}
\frac{\partial g_{\alpha}}{\partial y} & -\frac{\partial f_{\alpha}}{\partial y}\\
-\frac{\partial g_{\alpha}}{\partial x} & \frac{\partial f_{\alpha}}{\partial x}
\end{array}\right)\cdot\left(\begin{array}{c}
\frac{\partial\alpha}{\partial k}\\
\frac{\partial\alpha}{\partial k}
\end{array}\right)
\]
and therefore, by the assumptions $\frac{\partial g_{\alpha}}{\partial y}>\frac{\partial f_{\alpha}}{\partial y}\geq0$
and $\frac{\partial f_{\alpha}}{\partial x}>\frac{\partial g_{\alpha}}{\partial x}\geq0$
we have
\begin{align*}
\frac{\partial x}{\partial k} & =\frac{\frac{\partial\alpha}{\partial k}\cdot\left(\frac{\partial g_{\alpha}}{\partial y}-\frac{\partial f_{\alpha}}{\partial y}\right)}{\frac{\partial f_{\alpha}}{\partial x}\cdot\frac{\partial g_{\alpha}}{\partial y}-\frac{\partial f_{\alpha}}{\partial y}\cdot\frac{\partial g_{\alpha}}{\partial x}}>0\\
\frac{\partial y}{\partial k} & =\frac{\frac{\partial\alpha}{\partial k}\cdot\left(\frac{\partial f_{\alpha}}{\partial x}-\frac{\partial g_{\alpha}}{\partial x}\right)}{\frac{\partial f_{\alpha}}{\partial x}\cdot\frac{\partial g_{\alpha}}{\partial y}-\frac{\partial f_{\alpha}}{\partial y}\cdot\frac{\partial g_{\alpha}}{\partial x}}>0
\end{align*}
as required.
\end{proof}

\end{document}
